\documentclass[11pt]{amsart}

\usepackage[T1]{fontenc}
\usepackage{lmodern}
\usepackage{microtype}
\usepackage{mathtools,amssymb,amsthm}
\usepackage{array}
\usepackage[hidelinks]{hyperref}

\hypersetup{
  pdftitle={The cyclic antibandwidth of the 3x3x4 mesh is 13}
}

\newtheorem{theorem}{Theorem}
\newtheorem{lemma}[theorem]{Lemma}
\newtheorem{corollary}[theorem]{Corollary}
\newtheorem{proposition}[theorem]{Proposition}
\theoremstyle{remark}
\newtheorem*{remark}{Remark}

\emergencystretch=1.5em

\DeclareMathOperator{\cab}{cab}

\title[The cyclic antibandwidth of the $3\times3\times4$ mesh]
{The cyclic antibandwidth of the $3\times3\times4$ mesh is $13$}

\date{}

\subjclass[2020]{05C78, 05C85, 68R10}
\keywords{cyclic antibandwidth, antibandwidth, three-dimensional mesh, graph labeling,
exact value, exhaustive search}

\begin{document}

\begin{abstract}
For a graph $G$ on $n$ vertices, a labeling is a bijection $f\colon V(G)\to\{1,\dots,n\}$, and
the cyclic antibandwidth $\cab(G)$ is the largest $k$ for which some labeling places the two ends
of every edge at distance at least $k$ apart on the cycle $\mathbb Z_n$.
We prove $\cab(P_3\times P_3\times P_4)=13$. The lower bound is a printed $36$-label bijection,
certified by $75$ integer comparisons. The upper bound is a complete backtracking exhaustion at
$k=14$, using no solver and no library, run twice --- once under the symmetry reduction argued
here and once with no symmetry reduction at all; a degree count disposes of $k\ge16$ by hand.
The appendix table of conjectured $3$-dimensional-mesh values reproduced in the $2026$
exact-method paper of Xuan and Van lists $14$ at this instance, one above the value proved here.
What that table records, and by whom the entry was posed, are not settled here; nothing is
claimed about any other instance.
\end{abstract}

\maketitle

\section{Introduction}\label{sec:intro}

Let $G$ be a graph on $n=|V(G)|$ vertices. A \emph{labeling} of $G$ is a bijection
$f\colon V(G)\to\{1,\dots,n\}$. For labels $a,b$ put
\[
 d_c(a,b)=\min\bigl(|a-b|,\;n-|a-b|\bigr),
\]
the distance between $a$ and $b$ on the cycle $\mathbb Z_n$, and set
\[
 \cab_f(G)=\min_{uv\in E(G)}d_c\bigl(f(u),f(v)\bigr),
 \qquad
 \cab(G)=\max_f \cab_f(G).
\]
This is the \emph{cyclic antibandwidth} of $G$. It is a max-min parameter, so a single labeling
certifies a lower bound and nothing else: one edge sets the value of $\cab_f$, and a claimed
witness is checked in $O(|E|)$ integer arithmetic with no search. The $3$-dimensional mesh is the
Cartesian product $P_{n_1}\times P_{n_2}\times P_{n_3}$ of paths.

Exact values of $\cab$ are known for paths, cycles, two-dimensional and toroidal meshes and
hypercubes \cite{STV,RSSTV}, as surveyed in \cite{XT}; for three-dimensional meshes what the
literature offers is a table of conjectured values obtained from a heuristic. Nothing below uses
the contents of \cite{STV,RSSTV,TV}.

One such table is the source of the value we compare with. The exact-method paper of Xuan and Van
\cite{XT} reproduces, in its appendix on the experimental datasets, a $3$-dimensional mesh table
captioned ``Conjectured CAB value of 3D mesh instances. Source: Bansal et al.''; in the row for
$3\times3\times n_3$ meshes, read from the e-print \texttt{arXiv:2601.04239v1} (source lines
$1398$--$1422$), the entry at $n_3=4$ is $14$. In \cite{XT} the tabulated values are used to set
the search bound for $3$-dimensional meshes, which ``is chosen as the value immediately above the
conjectured CAB value reported by Bansal et al.''\ (source line $1395$), with the caveat that
``since this value is conjectured, it may be less than the actual optimal CAB value'' (line
$1534$).

\begin{theorem}\label{thm:main}
$\cab(P_3\times P_3\times P_4)=13$.
\end{theorem}

So the tabulated entry $14$ is one above the cyclic antibandwidth of that instance, and it is
unattained: no labeling of $P_3\times P_3\times P_4$ reaches $14$.

\begin{remark}
We do not claim to have refuted a stated conjecture. We did not read the body of \cite{BS}, which
is closed access, so we cannot say what closed form, if any, its authors posed; the accessible
source is the dataset caption quoted above. Nor do we establish which parameter the entry records.
The recorded run accompanying this note reports that the row's entries at the smallest six values
of $n_3$ coincide with values reported for the \emph{linear} antibandwidth of the same meshes, so
the printed $14$ may be a
linear rather than a cyclic value; on that reading it is not in conflict with anything proved here.
What is proved here is the displayed theorem.
\end{remark}

Section~\ref{sec:lower} exhibits a witness for $\cab\ge13$; six of its $75$ comparisons are
displayed and the rest are the same integer arithmetic.
Section~\ref{sec:upper} bounds $\cab$ above: a degree count gives $\cab\le15$ by hand, and a
complete backtracking exhaustion at $k=14$ of a reformulation (Lemma~\ref{lem:window}) gives
$\cab\le13$, using no solver, no library and no floating-point arithmetic.
Section~\ref{sec:status} says what is not settled.

\section{The lower bound: a printed witness}\label{sec:lower}

Write a vertex of $P_3\times P_3\times P_4$ as $(x,y,z)$ with $x,y\in\{1,2,3\}$ and
$z\in\{1,2,3,4\}$; two vertices are adjacent when they differ by $1$ in exactly one coordinate.
So $n=36$ and $|E|=75$. The labeling $W$ below is printed as four rows $z=1,\dots,4$; each row
groups $x=1\mid x=2\mid x=3$ and within a group lists $y=1,2,3$.

\begin{center}
$W$:\qquad
\begin{tabular}{r@{\;}|@{\;}ccc@{\;}|@{\;}ccc@{\;}|@{\;}ccc}
 $z=1$ & 31 & 14 & 30 & 13 & 29 & 10 & 35 & 12 & 27\\
 $z=2$ & 15 &  1 & 17 & 36 & 16 & 32 & 21 & 34 & 11\\
 $z=3$ &  2 & 23 &  4 & 22 &  3 & 18 &  7 & 20 & 33\\
 $z=4$ & 24 &  8 & 26 &  9 & 25 &  5 & 28 &  6 & 19
\end{tabular}
\end{center}

\begin{proposition}\label{prop:lower}
$W$ is a labeling of $P_3\times P_3\times P_4$ with $\cab_W=13$. Hence
$\cab(P_3\times P_3\times P_4)\ge13$.
\end{proposition}

\begin{proof}
The table lists each of $1,\dots,36$ exactly once, so $W$ is a bijection. For each of the $75$
edges one computes $d_c$ with $n=36$; the minimum is $13$, attained on $15$ edges. Nothing but
integer arithmetic is involved. As an illustration, take the
degree-$6$ vertex $(2,2,2)$, which carries label $16$; its six neighbours
\[
 (1,2,2),\;(3,2,2),\;(2,1,2),\;(2,3,2),\;(2,2,1),\;(2,2,3)
\]
carry $1,34,36,32,29,3$ respectively, and
\[
 d_c(16,\cdot)=15,\;18,\;16,\;16,\;13,\;13,
\]
all at least $13$.
\end{proof}

\section{The upper bound}\label{sec:upper}

Everything in this section rests on one elementary reformulation. For a label $l\in\mathbb Z_n$
put
\[
 F_k(l)=\{\,l'\in\mathbb Z_n : d_c(l,l')\ge k\,\}
      =\{\,l+k,\;l+k+1,\;\dots,\;l+n-k\,\},
\]
an \emph{arc} of exactly $n-2k+1$ labels antipodal to $l$.

\begin{lemma}\label{lem:window}
Let $G$ have $n$ vertices and let $1\le k\le\lfloor n/2\rfloor$. For a labeling $f$ the following
are equivalent.
\begin{enumerate}
\item[(i)] $\cab_f(G)\ge k$.
\item[(ii)] For every edge $uv$, $f(v)\in F_k(f(u))$; equivalently, $G$ is a spanning subgraph of
      the circulant $\mathrm{Ci}_n(k,k+1,\dots,\lfloor n/2\rfloor)$ under $f$.
\item[(iii)] For every $s\in\mathbb Z_n$ the set $f^{-1}(\{s,s+1,\dots,s+k-1\})$ of vertices
      carrying $k$ cyclically consecutive labels is independent in $G$.
\end{enumerate}
\end{lemma}

\begin{proof}
(i)$\Leftrightarrow$(ii) is the definition of $F_k$, and the circulant with connection set
$\{k,\dots,\lfloor n/2\rfloor\}$ is precisely the graph on $\mathbb Z_n$ whose edges join labels
at cyclic distance at least $k$. For (i)$\Leftrightarrow$(iii): two labels in a common window of
$k$ cyclically consecutive labels are at cyclic distance at most $k-1$, and conversely two labels
at cyclic distance $d\le k-1$ lie in a common such window. So some window carries an edge iff
some edge has $d_c<k$.
\end{proof}

Two consequences.

\begin{lemma}[$k\ge16$]\label{lem:cap}
$\cab(P_3\times P_3\times P_4)\le15$.
\end{lemma}

\begin{proof}
Let $\Delta=\Delta(G)$ and let $v$ have degree $\Delta$. By Lemma~\ref{lem:window}(ii) the $\Delta$
neighbours of $v$ receive distinct labels from the $n-2k+1$ labels of $F_k(f(v))$, so
$\cab(G)\le\lfloor(n-\Delta+1)/2\rfloor$. Here $n=36$ and $\Delta=6$, attained exactly at
$(2,2,2)$ and $(2,2,3)$, so $\cab\le\lfloor31/2\rfloor=15$. Concretely, $k=16$ leaves only
$36-31=5$ admissible labels for six neighbours.
\end{proof}

Neither of the two counting obstructions available here excludes $k=14$: the host circulant
$\mathrm{Ci}_{36}(14,\dots,18)$ is $9$-regular, which is at least $\Delta=6$, and by
Lemma~\ref{lem:window}(iii) a window of $14$ labels must be independent, which is no obstruction
either, since $P_3\times P_3\times P_4$ has a perfect matching and hence independence number
$36-18=18>14$. We decide $k=14$ by exhaustion.

\begin{lemma}[$k=14$]\label{lem:search}
There is no labeling $f$ of $P_3\times P_3\times P_4$ with $\cab_f\ge14$.
\end{lemma}

\begin{proof}[Proof (a complete backtracking exhaustion)]
By Lemma~\ref{lem:window}(ii) with $k=14$ we must decide whether $G=P_3\times P_3\times P_4$
embeds as a spanning subgraph of $\mathrm{Ci}_{36}(14,\dots,18)$. Order the $36$ vertices
$v_1,\dots,v_{36}$ deterministically: $v_1$ is the vertex of largest degree that is least in the
lexicographic order of $(x,y,z)$, so $v_1=(2,2,2)$; and $v_{i+1}$ is the unordered vertex with the
most already-ordered neighbours, ties broken by larger degree and then by lexicographic order.
Assign labels to $v_1,v_2,\dots$ in that order by depth-first search, trying admissible labels in
increasing order: a label is admissible for $v_i$ if it is unused and lies in $F_{14}(f(w))$ for
every already-labelled neighbour $w$ of $v_i$. After each assignment, propagate --- intersect the
candidate arc of every not-yet-labelled neighbour, and backtrack as soon as one of them is left
with no admissible label. This decides the question exactly: the search enumerates every
labeling with the required property, and it terminates.

Two reductions cut the tree, and both are exact.
\emph{Rotation.} $f\mapsto f+t$ preserves every cyclic distance and acts transitively on which
vertex carries a given label, so we may fix $f(v_1)=0$ (labels are taken in $\mathbb Z_{36}$
here).
\emph{The automorphisms of the layer.} The eight automorphisms of the $3\times3$ grid fix
$v_1=(2,2,2)$ and act transitively on its four in-layer neighbours, so we may require the least
of their four labels to sit on $(1,2,2)$; the one non-identity grid automorphism fixing $(1,2,2)$
transposes $(2,1,2)$ and $(2,3,2)$, so we may require $f(2,1,2)<f(2,3,2)$ as well. Finally
$f\mapsto-f$ fixes the label $0$, reverses every comparison and commutes with the grid
automorphisms, which do not touch $z$, so we may require $f(2,2,1)<f(2,2,3)$. Together these
quotient a group of order $16$ acting on the solution set.

Executed at $k=14$, the search reports no labeling in either run:
\[
\begin{array}{lrr}
 \text{reduced}, & 281{,}957{,}255 \text{ nodes}, & 0 \text{ labelings};\\
 \text{no symmetry reduction}, & 2{,}463{,}942{,}873 \text{ nodes}, & 0 \text{ labelings}.
\end{array}
\]
In the reduced run the tree never reaches beyond its $31$st level: every branch is
abandoned with at least five vertices still unlabelled. The second run uses no symmetry reduction
at all beyond the rotation pin, so it does not depend on the argument of the previous paragraph,
and it agrees.
\end{proof}

\begin{proof}[Proof of Theorem~\ref{thm:main}]
$\cab\ge13$ by Proposition~\ref{prop:lower}, and $\cab\le13$ by Lemma~\ref{lem:search}.
\end{proof}

\section{Status}\label{sec:status}

What is settled is $\cab(P_3\times P_3\times P_4)=13$. Every layer of that proof is re-derived by
the programs accompanying this note, from the labeling printed in Section~\ref{sec:lower}, with no
solver and no third-party library.

What is not settled:
\begin{itemize}
\item \emph{No other instance.} Nothing here bounds $\cab(P_3\times P_3\times P_{n_3})$ for any
      $n_3\ne4$, and the tabulated row is not repaired.
\item \emph{The status of the tabulated $14$.} Who posed it, in what form, and which parameter it
      records are all outside this note; see the remark in Section~\ref{sec:intro}.
\item \emph{The $k=14$ layer is a computation, not a readable proof.} It is solver-free,
      exhaustive and deterministic, but it is a search of hundreds of millions of nodes: a referee
      must re-run it rather than read it, and no independently checkable refutation object (a DRAT
      or LRAT trace, or a hand-sized infeasible core) accompanies it.
\end{itemize}

\section{Reproduction}

Everything up to and including Lemma~\ref{lem:cap} needs no code: rebuild
$P_3\times P_3\times P_4$, confirm $|V|=36$, $|E|=75$ and $\Delta=6$ at exactly $(2,2,2)$ and
$(2,2,3)$; confirm that the table of $W$ is a bijection onto $\{1,\dots,36\}$; compute the $75$
cyclic distances and read off the minimum $13$; and check the arithmetic of Lemma~\ref{lem:cap}.
That is one page, and it gives $13\le\cab\le15$.

The exhaustion of Lemma~\ref{lem:search} is machine work. Two programs accompany this note, each
with a recorded run: \texttt{verify.py} with \texttt{verify.output.txt} re-derives the objects and
the reduced $k=14$ exhaustion, and \texttt{exhaust.py} with \texttt{exhaust.output.txt} re-decides
$k=14$ with the symmetry reduction removed entirely. Both need Python~3.9 or later and the standard
library only --- no third-party package, no external data file, no network, no randomness and no
floating-point value in any decision --- and both are deterministic, so the node counts printed
above are reproducible exactly. Run them as \texttt{python3 verify.py} and
\texttt{python3 exhaust.py}; in the recorded runs they took about three minutes and about
twenty-three minutes on one core. Each program closes with a list of what it does not do; neither
checks a bibliographic locator, and neither bears on prior art.

\begin{thebibliography}{9}

\bibitem{BS}
R.~Bansal and K.~Srivastava,
\emph{A memetic algorithm for the cyclic antibandwidth maximization problem},
Soft Comput. \textbf{15} (2011), no.~2, 397--412.
\href{https://doi.org/10.1007/s00500-009-0538-6}{doi:10.1007/s00500-009-0538-6}.
Cited by \cite{XT} as the source of the tabulated values; the body is closed access and we did not
read it.

\bibitem{RSSTV}
A.~Raspaud, H.~Schr\"oder, O.~S\'ykora, L.~T\"or\"ok, and I.~Vr\v{t}o,
\emph{Antibandwidth and cyclic antibandwidth of meshes and hypercubes},
Discrete Math. \textbf{309} (2009), no.~11, 3541--3552.
\href{https://doi.org/10.1016/j.disc.2007.12.058}{doi:10.1016/j.disc.2007.12.058}.

\bibitem{STV}
O.~S\'ykora, L.~T\"or\"ok, and I.~Vr\v{t}o,
\emph{The cyclic antibandwidth problem},
Electron. Notes Discrete Math. \textbf{22} (2005), 223--227.
\href{https://doi.org/10.1016/j.endm.2005.06.030}{doi:10.1016/j.endm.2005.06.030}.

\bibitem{TV}
L.~T\"or\"ok and I.~Vr\v{t}o,
\emph{Antibandwidth of three-dimensional meshes},
Discrete Math. \textbf{310} (2010), no.~3, 505--510.
\href{https://doi.org/10.1016/j.disc.2009.03.029}{doi:10.1016/j.disc.2009.03.029}.

\bibitem{XT}
H.~T.~Xuan and K.~T.~Van,
\emph{Solving cyclic antibandwidth problem by SAT},
Comput. Optim. Appl. (2026),
\href{https://doi.org/10.1007/s10589-026-00819-8}{doi:10.1007/s10589-026-00819-8};
e-print \href{https://arxiv.org/abs/2601.04239}{arXiv:2601.04239v1}.
The caption and the entry quoted in Section~\ref{sec:intro} are read from the e-print source, in
the appendix table of conjectured $3$D-mesh CAB values, lines $1398$--$1422$.

\end{thebibliography}

\end{document}
